\section{Experiments}
\label{sec:experiments}

\subsection{Experimental protocol}

All numerical results reported in this section are generated by the project simulator and stored under \texttt{results/final\_2026-05-12/}. No synthetic or hand-adjusted data are introduced in the manuscript. The main experiment evaluates three holdout-fitted scenarios:
\begin{itemize}
    \item \texttt{paper\_base\_holdoutfit}: 120-slot episodes, synchronization budget 28, Eve speed noise standard deviation 0.8, and outage threshold 2.55;
    \item \texttt{paper\_hard\_holdoutfit}: 140-slot episodes, synchronization budget 22, Eve speed noise standard deviation 1.0, and outage threshold 2.72;
    \item \texttt{scenario\_stress\_holdoutfit}: 160-slot episodes, synchronization budget 19, Eve speed noise standard deviation 1.2, and calibrated outage threshold 1.10.
\end{itemize}
The main multi-seed table uses 20 seeds and compares \texttt{rollout\_joint}, \texttt{periodic}, \texttt{security\_risk}, and \texttt{security\_margin}. Additional stress-scenario studies use the same holdout-fitted stress configuration for SCA baselines, PPO, and strengthening ablations.

The primary metrics are average secrecy rate, secrecy outage probability, average synchronization cost, certificate cover rate, and runtime per slot. Higher secrecy rate is better, while lower outage and lower runtime are better. Confidence intervals are computed by the experiment scripts through bootstrap aggregation across seeds. Paired comparisons use matching random seeds for the compared methods.

\subsection{Main 20-seed results}

Table~\ref{tab:main_results} and Figure~\ref{fig:main_results} summarize the main 20-seed results. Across all three scenarios, \texttt{rollout\_joint} achieves the highest average secrecy rate among the four deployable methods. The gain is small in the base scenario, larger in the hard scenario, and strongest in the calibrated stress scenario.

\begin{table*}[t]
\centering
\caption{Main 20-seed results from \texttt{scheme\_c\_readiness\_20seed/main\_table.csv}.}
\label{tab:main_results}
\begin{tabular}{llrrrrr}
\toprule
Scenario & Method & Runs & Secrecy rate & Outage & Sync cost & Runtime ms/slot \\
\midrule
\multirow{4}{*}{Base} 
& rollout\_joint & 20 & 1.6429 & 0.8913 & 0.2018 & 1265.75 \\
& periodic & 20 & 1.6115 & 0.8913 & 0.1500 & 99.65 \\
& security\_risk & 20 & 1.4996 & 0.8896 & 0.0399 & 100.01 \\
& security\_margin & 20 & 0.6913 & 0.8967 & 0.2329 & 101.21 \\
\midrule
\multirow{4}{*}{Hard}
& rollout\_joint & 20 & 1.4429 & 0.9482 & 0.1032 & 1262.04 \\
& periodic & 20 & 1.2803 & 0.9496 & 0.1571 & 97.89 \\
& security\_risk & 20 & 1.2508 & 0.9432 & 0.0389 & 97.77 \\
& security\_margin & 20 & 0.4878 & 0.9657 & 0.1569 & 98.05 \\
\midrule
\multirow{4}{*}{Stress}
& rollout\_joint & 20 & 1.2540 & 0.4294 & 0.1098 & 1295.63 \\
& periodic & 20 & 0.9858 & 0.6434 & 0.1187 & 98.54 \\
& security\_risk & 20 & 0.8680 & 0.7003 & 0.0356 & 99.04 \\
& security\_margin & 20 & 0.6447 & 0.7856 & 0.1186 & 98.60 \\
\bottomrule
\end{tabular}
\end{table*}

\begin{figure*}[t]
\centering
\includegraphics[width=0.92\textwidth]{main_results_secrecy_outage.pdf}
\caption{Average secrecy rate and secrecy outage probability in the main 20-seed experiments. Error bars show the confidence intervals exported by the experiment scripts.}
\label{fig:main_results}
\end{figure*}

In the base scenario, \texttt{rollout\_joint} improves average secrecy rate over \texttt{periodic} by 0.0314, but both methods have the same outage probability. This indicates that the base scenario is not where the proposed method should be oversold; the fixed-period baseline is already strong under this operating condition. In the hard scenario, the secrecy-rate gain increases to 0.1626 with a small outage reduction of 0.0014. In the stress scenario, \texttt{rollout\_joint} increases average secrecy rate from 0.9858 to 1.2540 and reduces outage from 0.6434 to 0.4294 relative to \texttt{periodic}. The paired comparison file reports 20/20 secrecy-rate wins against \texttt{periodic}, \texttt{security\_risk}, and \texttt{security\_margin} in the stress scenario.

\subsection{Certificate validation}

Table~\ref{tab:certificate} and Figure~\ref{fig:certificate} report the holdout certificate validation results. The validation cover rates are above the target 0.95 level for all groups. The base validation group has cover rate 0.9873, while the hard and stress groups both reach 1.0000. The aggregate validation cover rate is 0.9964.

\begin{table}[t]
\centering
\caption{Holdout certificate validation from \texttt{scheme\_c\_holdout/holdout\_summary.csv}.}
\label{tab:certificate}
\begin{tabular}{llrr}
\toprule
Split & Scenario & Cover rate & Avg. realized loss \\
\midrule
validation & base & 0.9873 & 0.5075 \\
validation & hard & 1.0000 & 0.6568 \\
validation & stress & 1.0000 & 0.6916 \\
validation & all & 0.9964 & 0.6279 \\
\bottomrule
\end{tabular}
\end{table}

\begin{figure}[t]
\centering
\includegraphics[width=0.98\linewidth]{certificate_validation.pdf}
\caption{Holdout validation cover rates and average realized secrecy loss. The dashed line marks the 0.95 target cover rate.}
\label{fig:certificate}
\end{figure}

These results justify using the certificate as a conservative in-policy risk signal for the evaluated policies and scenarios. They do not justify claiming universal safety under arbitrary environment changes. If the Eve mobility model, channel model, or policy distribution changes substantially, the certificate should be refitted or revalidated.

\subsection{SCA and PPO baselines}

Table~\ref{tab:stress_baselines} reports additional stress-scenario baselines under the same holdout-fitted configuration. The SCA results show that local optimization is not sufficient in this setting. \texttt{sca\_oracle} obtains secrecy rate 1.2126, close to \texttt{rollout\_joint}'s 1.2558, but its outage probability is much worse (0.8125 versus 0.4200). This indicates that improving average secrecy alone does not guarantee reliable secure service. The deployable \texttt{sca\_twin} baseline performs poorly, with secrecy rate 0.4677 and outage 0.8225.

\begin{table*}[t]
\centering
\caption{Stress holdout-fitted SCA and PPO baselines. Values are from the corresponding \texttt{main\_table.csv} files.}
\label{tab:stress_baselines}
\begin{tabular}{llrrrr}
\toprule
Group & Method & Runs & Secrecy rate & Outage & Runtime ms/slot \\
\midrule
\multirow{4}{*}{SCA}
& rollout\_joint & 5 & 1.2558 & 0.4200 & 1316.34 \\
& sca\_oracle & 5 & 1.2126 & 0.8125 & 246.50 \\
& periodic & 5 & 0.9955 & 0.6563 & 99.63 \\
& sca\_twin & 5 & 0.4677 & 0.8225 & 252.27 \\
\midrule
\multirow{3}{*}{PPO}
& rollout\_joint & 5 & 1.2558 & 0.4200 & 1318.07 \\
& periodic & 5 & 0.9955 & 0.6563 & 98.85 \\
& ppo\_baseline & 5 & 0.3442 & 0.8463 & 6.33 \\
\bottomrule
\end{tabular}
\end{table*}

The PPO baseline is very fast at inference time (6.33 ms/slot), but its secrecy rate is only 0.3442 under the 200-episode training budget. Figure~\ref{fig:ppo_training} shows the training trajectory exported by the PPO script. The curve supports a conservative interpretation: the implemented lightweight PPO baseline did not learn a competitive policy within the given budget, but this should not be read as a general rejection of deep reinforcement learning for the problem.

\begin{figure}[t]
\centering
\includegraphics[width=0.98\linewidth]{ppo_training_curve.pdf}
\caption{PPO training curve from \texttt{drl\_ppo\_200ep\_5seed\_stress\_holdoutfit/training\_log.csv}. The smoothed curve shows the 10-episode moving average.}
\label{fig:ppo_training}
\end{figure}

\subsection{Strengthening and ablation suite}

Table~\ref{tab:ablation} and Figure~\ref{fig:strengthening} report the stress holdout-fitted strengthening suite. The oracle-state rollout achieves the best result, with secrecy rate 1.3121 and outage 0.0292, and is treated as an upper reference rather than a deployable method. Among deployable methods, \texttt{rollout\_joint} is strongest. It improves over \texttt{rollout\_fixed\_periodic} by 0.0408 in secrecy rate and reduces outage from 0.5396 to 0.4083, showing that adaptive synchronization contributes beyond rollout motion and power search. It also substantially outperforms \texttt{no\_twin} and \texttt{rollout\_no\_sync}, confirming that twin prediction and synchronization are both important.

\begin{table*}[t]
\centering
\caption{Stress holdout-fitted strengthening suite from \texttt{strengthening\_suite\_3seed\_stress\_holdoutfit/main\_table.csv}.}
\label{tab:ablation}
\begin{tabular}{lrrrr}
\toprule
Method & Runs & Secrecy rate & Outage & Runtime ms/slot \\
\midrule
oracle\_sync & 3 & 1.3121 & 0.0292 & 1412.94 \\
rollout\_joint & 3 & 1.2550 & 0.4083 & 1305.39 \\
rollout\_fixed\_periodic & 3 & 1.2142 & 0.5396 & 1208.12 \\
periodic & 3 & 1.0210 & 0.6417 & 98.84 \\
no\_twin & 3 & 0.9597 & 0.7021 & 1247.44 \\
rollout\_no\_sync & 3 & 0.6670 & 0.8229 & 775.95 \\
sca\_twin & 3 & 0.4708 & 0.8208 & 247.08 \\
random\_budgeted & 3 & 0.0046 & 1.0000 & 6.30 \\
\bottomrule
\end{tabular}
\end{table*}

\begin{figure*}[t]
\centering
\includegraphics[width=0.96\textwidth]{stress_strengthening_suite.pdf}
\caption{Stress-scenario strengthening suite. The proposed method is below the oracle-state upper reference but outperforms the deployable ablations in secrecy rate and outage.}
\label{fig:strengthening}
\end{figure*}

\subsection{Small exact-DP sanity check}

To verify the dynamic-programming evaluation chain, we also run a deliberately small exact-DP instance. The configuration evaluates 56,784 abstract states and 72 actions over an episode length of 4, obtaining optimal cumulative secrecy 9.5134 and optimal average secrecy 2.3783. Figure~\ref{fig:small_mdp} shows the resulting optimal trace. This experiment is not a paper-scale upper bound because its horizon, action space, and abstraction differ from the main scenarios. Its purpose is to confirm that the exact-DP machinery works on a small instance.

\begin{figure}[t]
\centering
\includegraphics[width=0.98\linewidth]{small_mdp_trace.pdf}
\caption{Small exact-DP sanity-check trace from \texttt{small\_mdp\_bound/optimal\_trace.csv}.}
\label{fig:small_mdp}
\end{figure}

\subsection{Runtime tradeoff}

The main limitation of the proposed method is runtime. In the main 20-seed table, \texttt{rollout\_joint} requires approximately 1265.75, 1262.04, and 1295.63 ms/slot in the base, hard, and stress scenarios, respectively. The corresponding \texttt{periodic} runtimes are about 99.65, 97.89, and 98.54 ms/slot. Thus, the proposed method is better interpreted as a performance-oriented controller for edge-assisted or moderate-timescale decision making, rather than as an ultra-low-latency onboard controller in its current implementation.
